\chapter{Por onde começar}
\markboth{Começar}{Começar}

\section{Montar e rodar}

\begin{verbatim}
npm install
make -C native test        # suite de aceitacao herdada: a rede de seguranca
npm run wasm               # precisa do emsdk: source ~/emsdk/emsdk_env.sh
npm test && npm run build

./native/build/bin/ghost run specs/neutrino-vacuo.json --saida saida.csv
\end{verbatim}

Rode \cmd{make -C native test} antes de qualquer outra coisa, e rode de novo
depois de cada mudança. É a suíte herdada do DAEDALUS e é o que impede que uma
alteração na superfície quebre o núcleo em silêncio.

\begin{nota}
\rotulo{o spec.json é a moeda única}
Um cálculo é descrito por um \cmd{spec.json}, e ele é analisado \emph{só em C}.
Nada de montar hamiltoniano em Python: um script de figuras escreve o spec,
chama o binário e desenha. Ele não calcula física.

\medskip
O cabeçalho de todo CSV carrega o \emph{hash do núcleo}. Um resultado do GHOST
não pode ser confundido com um do DAEDALUS, nem reimportado como se fosse. É
para isso que o campo existe: proveniência é conteúdo, não decoração.
\end{nota}

\section{Uma sequência de exercícios}

Em ordem. Cada um usa o anterior.

\begin{enumerate}
\item \textbf{Dois sabores no vácuo.} Reproduza
  $P = \sin^2 2\theta \sin^2(kL)$ e confirme o erro na casa $10^{-15}$. Se der
  $10^{-8}$, algo está em precisão simples.

\item \textbf{Varra o ângulo.} Refaça a figura da página~\pageref{cap:teoria}:
  a profundidade do vale contra $\sin^2 2\theta$. Convença-se de que matéria
  \emph{não} muda a mistura.

\item \textbf{Ligue a matéria, densidade constante.} Encontre a ressonância
  varrendo $\rho$ e verifique a equação~\eqref{eq:resonancia}. Veja a amplitude
  ir de $0{,}117$ a $0{,}999$ com $\theta$ fixo.

\item \textbf{Perfil exponencial.} Compare contra a fórmula de Petcov. Aqui
  aparece pela primeira vez a questão do domínio de integração --- e, se você
  errar o domínio, aparece também o erro de 26\%.

\item \textbf{Três sabores com fase CP.} Passa a usar o mergulho real
  $2n\times2n$. Compare contra \cmd{expm}.

\item \textbf{Trajetórias.} Ligue a defasagem e recupere a solução exata de
  Lindblad. Confira o traço: se não der $1{,}000000000000$, pare.
\end{enumerate}

\begin{alerta}
\rotuloa{o erro que você vai cometer no exercício 4}
Convergir olhando $P(\nue)$. Ela oscila e engana --- com 1500 fatias deu
$0{,}0027$ contra os $0{,}0224$ corretos. A métrica é
$|\braket{\nu_3}{\psi}|^2 \to 1$. Está no capítulo~\ref{cap:nucleo} e está
repetido aqui porque é o erro mais caro do projeto.
\end{alerta}

\section{Bibliografia comentada}

\subsection{Para entender o efeito}

\begin{description}[leftmargin=0pt, itemsep=0.7ex]
\item[Wolfenstein, 1978.] O potencial de matéria. Curto. Leia o argumento de
  por que a corrente neutra não conta.
\item[Mikheyev \& Smirnov, 1985.] A ressonância. O ponto não é a fórmula, é
  perceber que densidade \emph{variável} é um problema diferente de densidade
  constante.
\item[Petcov, 1988; Kuo \& Pantaleone, 1989.] A fórmula fechada para o perfil
  exponencial. É contra ela que o código se valida no exercício 4.
\end{description}

\subsection{Para entender o choque}

\begin{description}[leftmargin=0pt, itemsep=0.7ex]
\item[Schirato \& Fuller, 2002.] O choque quebra a adiabaticidade. O artigo
  fundador desta linha.
\item[Fogli et al., 2003.] A assinatura migra em energia: baixa primeiro,
  depois alta. É o mapa em $(t,E)$, e é deles.
\item[Tomàs et al., 2004.] Choque reverso, estrutura não monótona no tempo, e a
  primeira menção de que a frente numérica precisa ser afiada.
\item[Barger, Huber \& Marfatia, 2005.] Enunciam a prática do afiamento. Não
  quantificam a ausência dela --- e é aí que fica o espaço.
\item[Kneller \& Volpe, 2010; Kneller \& Mauney, 2013.] Os limiares de amplitude
  para turbulência. Leia antes de propor qualquer coisa sobre turbulência.
\end{description}

\subsection{Para o setor escuro e NSI}

Ver a seção~\ref{sec:setor-escuro}. Esta lista existe para evitar que a
direção seja reproposta como se fosse nova.

\begin{description}[leftmargin=0pt, itemsep=0.7ex]
\item[Esteban-Pretel, Tomás \& Valle, 2007.] \emph{Probing non-standard
  neutrino interactions with supernova neutrinos}, arXiv:0704.0032. \textbf{Leia
  este primeiro.} NSI na propagação pelo envelope, conversão ressonante induzida
  por NSI, e a passagem da onda de choque --- para NSI não-universais de poucos
  por cento. É o trabalho que já ocupa o efeito.
\item[Esteban-Pretel, Tomás \& Valle, 2009.] arXiv:0909.2196. O mesmo, em
  interação com os efeitos coletivos.
\item[Jana \& Porto, 2024.] \emph{Non-standard interactions of supernova
  neutrinos and mass ordering ambiguity at DUNE}, arXiv:2407.06251. NSI inverte
  os níveis de matéria e induz nova ressonância perto da protoestrela; é de onde
  vem a sensibilidade $\mathcal{O}(0{,}1)$ marcada na
  figura~\ref{fig:eps-nsi}.
\item[Dutta \& Karthikeyan, 2025 (arXiv:2508.16558); Das \& Reno, 2025
  (arXiv:2508.16510).] A variante escalar, por dois grupos no mesmo mês. Sinal
  de área quente.
\item[Sen, 2024.] \emph{Supernova Neutrinos: Flavour Conversion Mechanisms and
  New Physics Scenarios}, arXiv:2405.20432. Revisão; o mapa do que está
  ocupado.
\item[Volpe, 2023.] \emph{Neutrinos from dense environments}, arXiv:2301.11814.
  Revisão mais ampla, de mecanismos e abordagens teóricas.
\item[Chang, Essig \& McDermott, 2016 (arXiv:1611.03864) e 2018
  (arXiv:1803.00993).] Os vínculos de resfriamento da SN 1987A para fótons
  escuros e setores escuros sub-GeV. É a via de resfriamento, feita com a
  ferramenta certa.
\item[Akita, Im, Masud \& Yun, 2023.] \emph{Limits on heavy neutral leptons,
  $Z'$ bosons and majorons from high-energy supernova neutrinos},
  arXiv:2312.13627. A via de decaimento, feita, e com exatamente este conteúdo
  de partículas.
\item[Farzan, 2015.] arXiv:1505.06906. Modelo de mediador leve que gera NSI
  grande; a ponte entre um $Z'$ e o parâmetro $\varepsilon$.
\item[Acero et al., 2022.] \emph{White paper on light sterile neutrino
  searches}, arXiv:2203.07323. Onde o cenário HNL mais fóton escuro do
  MiniBooNE está catalogado, com o espaço de parâmetros.
\end{description}

\subsection{Para o lado numérico}

\begin{description}[leftmargin=0pt, itemsep=0.7ex]
\item[\cmd{CONVENTIONS.md}, no próprio repositório.] Não é documentação de
  apoio; é a regra. Um único núcleo, spec analisado só em C,
  \cmd{-ffast-math} proibido, sem \cmd{\_Complex}, zero estado global,
  reprodutibilidade bit a bit, limite espectral garantido.
\end{description}

\subsection{Os dados}

\begin{description}[leftmargin=0pt, itemsep=0.7ex]
\item[Garching Core-Collapse Supernova Archive.] Modelos 3D de
  A.~Wongwathanarat em grade Yin-Yang. \textbf{Atenção:} o que é distribuído
  publicamente são perfis promediados em ângulo, e a média destrói a frente por
  um fator 18. Para este trabalho é preciso o raio individual.
\end{description}

\section{Uma última coisa}

\begin{conceito}
\rotuloc{se você lembrar de três coisas}
\textbf{Uma.} O neutrino tem dois conjuntos de estados e eles discordam; tudo o
mais é consequência.

\smallskip
\textbf{Duas.} Na ressonância, o que decide não é o ângulo de mistura, é a
velocidade com que a densidade muda.

\smallskip
\textbf{Três.} O observável é o incoerente, e a métrica de convergência é
$|\braket{\nu_3}{\psi}|^2$. As duas convenções que a física impôs ao código
existem porque a alternativa produz números errados sem avisar.
\end{conceito}
